\documentclass[11pt,a4paper]{article}

%=========================================================
% PACKAGES
%=========================================================
\usepackage[utf8]{inputenc}
\usepackage[T1]{fontenc}
\usepackage[french]{babel}
\usepackage[a4paper,margin=2.15cm]{geometry}
\usepackage{lmodern}
\usepackage{amsmath,amssymb,amsthm,mathtools}
\usepackage{fancyhdr}
\usepackage{lastpage}
\usepackage{enumitem}
\usepackage{array,tabularx}
\usepackage{tikz}
\usepackage[most]{tcolorbox}
\usepackage{hyperref}
\usepackage{xcolor}

%=========================================================
% MISE EN PAGE
%=========================================================
\setlength{\headheight}{14pt}
\pagestyle{fancy}
\fancyhf{}
\fancyhead[L]{Analyse CPGE}
\fancyhead[C]{Suites numériques}
\fancyhead[R]{Page \thepage/\pageref{LastPage}}
\fancyfoot[C]{Convergence -- Monotonie -- Extractions -- Comparaisons}

\hypersetup{
    colorlinks=true,
    linkcolor=blue!50!black,
    urlcolor=blue!50!black
}

\setlist[itemize]{label=\(\triangleright\)}
\setlist[enumerate]{label=\textbf{\arabic*.}}

%=========================================================
% COULEURS ET BOÎTES
%=========================================================
\definecolor{bleufonce}{RGB}{25,65,130}
\definecolor{bleuclair}{RGB}{235,243,255}
\definecolor{vertfonce}{RGB}{30,120,70}
\definecolor{vertclair}{RGB}{235,250,240}
\definecolor{rougefonce}{RGB}{170,40,40}
\definecolor{rougeclair}{RGB}{255,238,238}
\definecolor{orangefonce}{RGB}{190,110,20}
\definecolor{orangeclair}{RGB}{255,246,230}
\definecolor{violetfonce}{RGB}{95,55,145}
\definecolor{violetclair}{RGB}{246,241,255}

\tcbset{
    enhanced,
    sharp corners,
    boxrule=0.75pt,
    before skip=8pt,
    after skip=8pt
}

\newtcolorbox{definitionbox}[1][]{
    colback=bleuclair,
    colframe=bleufonce,
    title=\textbf{Définition #1}
}

\newtcolorbox{theorembox}[1][]{
    colback=vertclair,
    colframe=vertfonce,
    title=\textbf{Théorème / Propriété #1}
}

\newtcolorbox{methodbox}[1][]{
    colback=orangeclair,
    colframe=orangefonce,
    title=\textbf{Méthode #1}
}

\newtcolorbox{retenirbox}[1][]{
    colback=violetclair,
    colframe=violetfonce,
    title=\textbf{À retenir #1}
}

\newtcolorbox{erreurbox}[1][]{
    colback=rougeclair,
    colframe=rougefonce,
    title=\textbf{Erreur classique #1}
}

\newtcolorbox{exercicebox}[1][]{
    colback=white,
    colframe=black!60,
    title=\textbf{Exercice #1}
}

\newtcolorbox{correctionbox}[1][]{
    colback=gray!5,
    colframe=black!65,
    title=\textbf{Correction #1}
}

%=========================================================
% COMMANDES
%=========================================================
\newcommand{\N}{\mathbb N}
\newcommand{\Z}{\mathbb Z}
\newcommand{\Q}{\mathbb Q}
\newcommand{\R}{\mathbb R}
\newcommand{\C}{\mathbb C}
\newcommand{\abs}[1]{\left|#1\right|}
\newcommand{\norme}[1]{\left\lVert #1\right\rVert}
\newcommand{\ssi}{\Longleftrightarrow}
\newcommand{\implique}{\Longrightarrow}
\newcommand{\limn}{\lim_{n\to+\infty}}

\newtheorem{prop}{Proposition}[section]
\newtheorem{lemme}[prop]{Lemme}
\newtheorem{corollaire}[prop]{Corollaire}

%=========================================================
% DOCUMENT
%=========================================================
\begin{document}

\begin{center}
    {\Huge \textbf{Suites numériques}}\\[2mm]
    {\Large Convergence, monotonie, suites extraites et comparaisons}\\[2mm]
    {\large Cours complet pour classes préparatoires scientifiques}\\[2mm]
    \rule{0.88\linewidth}{0.5pt}
\end{center}

\vspace{0.4cm}

\tableofcontents

\newpage

%=========================================================
\section*{Introduction générale}
\addcontentsline{toc}{section}{Introduction générale}

Les suites numériques constituent le premier grand chapitre d'analyse. Elles permettent de formaliser l'idée intuitive suivante :

\begin{center}
\textit{Que devient une quantité lorsque l'indice devient très grand ?}
\end{center}

Une suite est une famille de nombres indexée par les entiers naturels :
\[
u_0,u_1,u_2,\dots,u_n,\dots
\]
L'objectif principal est d'étudier son comportement lorsque :
\[
n\to+\infty.
\]

Ce chapitre introduit les notions fondamentales qui seront reprises dans tout le reste de l'analyse :
\begin{itemize}
    \item convergence ;
    \item limite ;
    \item unicité de la limite ;
    \item opérations sur les limites ;
    \item encadrement ;
    \item monotonie ;
    \item suites adjacentes ;
    \item suites extraites ;
    \item théorème de Bolzano-Weierstrass ;
    \item suites récurrentes ;
    \item comparaisons asymptotiques.
\end{itemize}

\begin{retenirbox}
Les suites sont le laboratoire de la notion de limite.  
Avant de manipuler des limites de fonctions, des séries ou des intégrales, il faut maîtriser les suites.
\end{retenirbox}

%=========================================================
\section{Définitions fondamentales}

\subsection{Suite réelle ou complexe}

\begin{definitionbox}[Suite numérique]
Une suite numérique est une application :
\[
u:\N\to\R
\]
ou :
\[
u:\N\to\C.
\]
On note généralement :
\[
u_n=u(n).
\]
La suite est alors notée :
\[
(u_n)_{n\in\N}.
\]
\end{definitionbox}

Si les termes appartiennent à \(\R\), on parle de suite réelle.  
Si les termes appartiennent à \(\C\), on parle de suite complexe.

\begin{retenirbox}
Une suite n'est pas seulement une liste de nombres : c'est une fonction définie sur \(\N\).
\[
n\longmapsto u_n.
\]
\end{retenirbox}

\subsection{Exemples classiques}

\[
u_n=\frac1n,\quad n\geq1.
\]
Cette suite tend vers \(0\).

\[
v_n=(-1)^n.
\]
Cette suite alterne entre \(1\) et \(-1\), donc elle ne converge pas.

\[
w_n=n^2.
\]
Cette suite tend vers \(+\infty\).

\[
z_n=\left(\frac{1+i}{2}\right)^n.
\]
Cette suite est complexe.

\subsection{Suite définie explicitement et suite définie par récurrence}

\begin{definitionbox}[Suite explicite]
Une suite est définie explicitement lorsque \(u_n\) est donné directement en fonction de \(n\).
\[
u_n=\frac{n}{n+1}.
\]
\end{definitionbox}

\begin{definitionbox}[Suite récurrente]
Une suite est définie par récurrence lorsqu'un premier terme est donné, puis une relation permet de calculer les termes suivants :
\[
u_0\in\R,\qquad u_{n+1}=f(u_n).
\]
\end{definitionbox}

\begin{erreurbox}
Dans une suite récurrente, on ne connaît pas directement \(u_n\).  
Il faut souvent commencer par étudier :
\[
u_{n+1}-u_n,\qquad f(x)-x,\qquad \text{ou les points fixes de }f.
\]
\end{erreurbox}

%=========================================================
\section{Suites bornées, majorées et minorées}

\subsection{Définitions}

\begin{definitionbox}[Suite majorée]
Une suite réelle \((u_n)\) est majorée s'il existe \(M\in\R\) tel que :
\[
\forall n\in\N,\qquad u_n\leq M.
\]
\end{definitionbox}

\begin{definitionbox}[Suite minorée]
Une suite réelle \((u_n)\) est minorée s'il existe \(m\in\R\) tel que :
\[
\forall n\in\N,\qquad m\leq u_n.
\]
\end{definitionbox}

\begin{definitionbox}[Suite bornée]
Une suite réelle \((u_n)\) est bornée si elle est à la fois majorée et minorée.
\end{definitionbox}

\begin{theorembox}[Caractérisation par la valeur absolue]
Une suite réelle \((u_n)\) est bornée si et seulement s'il existe \(M\geq0\) tel que :
\[
\forall n\in\N,\qquad \abs{u_n}\leq M.
\]
\end{theorembox}

\begin{proof}
Supposons que \((u_n)\) soit bornée.  
Il existe \(m,M\in\R\) tels que :
\[
\forall n\in\N,\qquad m\leq u_n\leq M.
\]
Posons :
\[
R=\max(\abs{m},\abs{M}).
\]
Alors :
\[
\forall n\in\N,\qquad \abs{u_n}\leq R.
\]

Réciproquement, s'il existe \(R\geq0\) tel que :
\[
\abs{u_n}\leq R,
\]
alors :
\[
-R\leq u_n\leq R.
\]
Donc la suite est minorée par \(-R\) et majorée par \(R\).
\end{proof}

\begin{exercicebox}[Suite bornée]
Montrer que la suite :
\[
u_n=\frac{n}{n+1}
\]
est bornée.
\end{exercicebox}

\begin{correctionbox}
Pour tout \(n\in\N\), on a :
\[
0\leq \frac{n}{n+1}<1.
\]
Donc :
\[
0\leq u_n\leq1.
\]
La suite est donc minorée par \(0\) et majorée par \(1\).  
Elle est bornée.
\end{correctionbox}

%=========================================================
\section{Convergence d'une suite réelle}

\subsection{Définition de la convergence}

\begin{definitionbox}[Convergence vers une limite réelle]
Soit \((u_n)\) une suite réelle et soit \(\ell\in\R\).  
On dit que \((u_n)\) converge vers \(\ell\) si :
\[
\forall \varepsilon>0,\quad \exists N\in\N,\quad \forall n\geq N,\quad \abs{u_n-\ell}\leq \varepsilon.
\]
On note :
\[
u_n\longrightarrow \ell
\]
ou :
\[
\lim_{n\to+\infty}u_n=\ell.
\]
\end{definitionbox}

\begin{retenirbox}
La définition signifie que les termes de la suite finissent par être arbitrairement proches de \(\ell\).
\[
\abs{u_n-\ell}\leq\varepsilon
\]
signifie que \(u_n\) appartient à l'intervalle :
\[
[\ell-\varepsilon,\ell+\varepsilon].
\]
\end{retenirbox}

\subsection{Interprétation de la définition}

La phrase :
\[
\forall \varepsilon>0,\quad \exists N\in\N,\quad \forall n\geq N,\quad \abs{u_n-\ell}\leq \varepsilon
\]
se lit :

\begin{center}
Pour toute marge d'erreur \(\varepsilon>0\), aussi petite soit-elle, il existe un rang \(N\) à partir duquel tous les termes de la suite sont à distance au plus \(\varepsilon\) de \(\ell\).
\end{center}

\begin{erreurbox}
Le rang \(N\) peut dépendre de \(\varepsilon\).  
On ne demande pas qu'un même \(N\) fonctionne pour tous les \(\varepsilon>0\).
\end{erreurbox}

\subsection{Exemple détaillé}

\begin{exercicebox}[Définition de la limite]
Montrer avec la définition que :
\[
\frac1n\longrightarrow0.
\]
\end{exercicebox}

\begin{correctionbox}
On veut montrer :
\[
\forall \varepsilon>0,\quad \exists N\in\N,\quad \forall n\geq N,\quad \left|\frac1n-0\right|\leq\varepsilon.
\]
Soit \(\varepsilon>0\).  
On cherche \(N\) tel que :
\[
\frac1n\leq\varepsilon
\]
pour tout \(n\geq N\).  
Il suffit d'avoir :
\[
n\geq \frac1\varepsilon.
\]
On choisit donc un entier \(N\) tel que :
\[
N\geq \frac1\varepsilon.
\]
Alors, pour tout \(n\geq N\),
\[
0<\frac1n\leq\frac1N\leq\varepsilon.
\]
Donc :
\[
\frac1n\to0.
\]
\end{correctionbox}

\subsection{Suite convergente et suite divergente}

\begin{definitionbox}[Suite convergente]
Une suite est convergente s'il existe \(\ell\in\R\) tel que :
\[
u_n\to\ell.
\]
\end{definitionbox}

\begin{definitionbox}[Suite divergente]
Une suite qui n'est pas convergente est dite divergente.
\end{definitionbox}

\begin{erreurbox}
Dire qu'une suite diverge ne signifie pas forcément qu'elle tend vers \(+\infty\) ou \(-\infty\).  
Par exemple :
\[
(-1)^n
\]
diverge, mais ne tend ni vers \(+\infty\), ni vers \(-\infty\).
\end{erreurbox}

%=========================================================
\section{Unicité de la limite et suites convergentes bornées}

\subsection{Unicité de la limite}

\begin{theorembox}[Unicité de la limite]
Si une suite réelle converge, alors sa limite est unique.
\end{theorembox}

\begin{proof}
Supposons que :
\[
u_n\to \ell
\]
et :
\[
u_n\to m.
\]
Montrons que :
\[
\ell=m.
\]

Soit \(\varepsilon>0\).  
Comme \(u_n\to\ell\), il existe \(N_1\in\N\) tel que :
\[
n\geq N_1\implique \abs{u_n-\ell}\leq \frac{\varepsilon}{2}.
\]
Comme \(u_n\to m\), il existe \(N_2\in\N\) tel que :
\[
n\geq N_2\implique \abs{u_n-m}\leq \frac{\varepsilon}{2}.
\]
Pour \(n\geq \max(N_1,N_2)\), on a :
\[
\abs{\ell-m}
=
\abs{\ell-u_n+u_n-m}.
\]
Par inégalité triangulaire :
\[
\abs{\ell-m}
\leq
\abs{\ell-u_n}+\abs{u_n-m}
\leq
\frac{\varepsilon}{2}+\frac{\varepsilon}{2}
=
\varepsilon.
\]
Donc :
\[
\forall \varepsilon>0,\qquad \abs{\ell-m}\leq\varepsilon.
\]
Cela impose :
\[
\abs{\ell-m}=0.
\]
Donc :
\[
\ell=m.
\]
\end{proof}

\subsection{Toute suite convergente est bornée}

\begin{theorembox}
Toute suite réelle convergente est bornée.
\end{theorembox}

\begin{proof}
Supposons :
\[
u_n\to\ell.
\]
En prenant \(\varepsilon=1\), il existe \(N\in\N\) tel que :
\[
\forall n\geq N,\qquad \abs{u_n-\ell}\leq1.
\]
Alors, pour \(n\geq N\),
\[
\abs{u_n}\leq\abs{\ell}+1.
\]
Les termes :
\[
u_0,u_1,\dots,u_{N-1}
\]
sont en nombre fini, donc leurs valeurs absolues admettent un maximum :
\[
M_0=\max(\abs{u_0},\dots,\abs{u_{N-1}}).
\]
Posons :
\[
M=\max(M_0,\abs{\ell}+1).
\]
Alors :
\[
\forall n\in\N,\qquad \abs{u_n}\leq M.
\]
Donc la suite est bornée.
\end{proof}

\begin{erreurbox}
La réciproque est fausse : une suite bornée n'est pas forcément convergente.  
Exemple :
\[
u_n=(-1)^n.
\]
Elle est bornée, mais elle ne converge pas.
\end{erreurbox}

%=========================================================
\section{Limites infinies}

\subsection{Définition}

\begin{definitionbox}[Tendre vers \(+\infty\)]
Une suite réelle \((u_n)\) tend vers \(+\infty\) si :
\[
\forall A\in\R,\quad \exists N\in\N,\quad \forall n\geq N,\quad u_n\geq A.
\]
On note :
\[
u_n\to+\infty.
\]
\end{definitionbox}

\begin{definitionbox}[Tendre vers \(-\infty\)]
Une suite réelle \((u_n)\) tend vers \(-\infty\) si :
\[
\forall A\in\R,\quad \exists N\in\N,\quad \forall n\geq N,\quad u_n\leq A.
\]
On note :
\[
u_n\to-\infty.
\]
\end{definitionbox}

\begin{retenirbox}
Dire :
\[
u_n\to+\infty
\]
ne signifie pas que \(+\infty\) est une limite réelle.  
Cela signifie que les termes deviennent arbitrairement grands.
\end{retenirbox}

\begin{exercicebox}[Limite infinie]
Montrer que :
\[
n^2\to+\infty.
\]
\end{exercicebox}

\begin{correctionbox}
Soit \(A\in\R\).  
On cherche \(N\) tel que :
\[
n^2\geq A
\]
pour tout \(n\geq N\).

Si \(A\leq0\), tout \(N\) convient, par exemple \(N=0\).  
Si \(A>0\), il suffit de prendre :
\[
N\geq \sqrt A.
\]
Alors, pour tout \(n\geq N\),
\[
n^2\geq N^2\geq A.
\]
Donc :
\[
n^2\to+\infty.
\]
\end{correctionbox}

%=========================================================
\section{Opérations sur les limites}

\subsection{Somme, produit, quotient}

\begin{theorembox}[Opérations sur les limites finies]
Soient \((u_n)\) et \((v_n)\) deux suites réelles telles que :
\[
u_n\to \ell
\quad\text{et}\quad
v_n\to m.
\]
Alors :
\[
u_n+v_n\to \ell+m,
\]
\[
u_nv_n\to \ell m.
\]
Si \(m\neq0\) et si \(v_n\neq0\) à partir d'un certain rang, alors :
\[
\frac{u_n}{v_n}\to \frac{\ell}{m}.
\]
\end{theorembox}

\begin{proof}
Pour la somme :
\[
\abs{(u_n+v_n)-(\ell+m)}
\leq
\abs{u_n-\ell}+\abs{v_n-m}.
\]
Soit \(\varepsilon>0\).  
À partir d'un certain rang :
\[
\abs{u_n-\ell}\leq\frac{\varepsilon}{2}
\]
et :
\[
\abs{v_n-m}\leq\frac{\varepsilon}{2}.
\]
Donc :
\[
\abs{(u_n+v_n)-(\ell+m)}\leq\varepsilon.
\]

Pour le produit, on écrit :
\[
u_nv_n-\ell m
=
u_n(v_n-m)+m(u_n-\ell).
\]
Donc :
\[
\abs{u_nv_n-\ell m}
\leq
\abs{u_n}\abs{v_n-m}+\abs{m}\abs{u_n-\ell}.
\]
Comme \((u_n)\) est convergente, elle est bornée.  
Il existe donc \(M\geq0\) tel que :
\[
\abs{u_n}\leq M.
\]
Les deux termes du membre de droite tendent vers \(0\), donc :
\[
u_nv_n\to\ell m.
\]

Pour le quotient, on utilise :
\[
\frac{u_n}{v_n}=u_n\cdot\frac1{v_n}.
\]
Il suffit donc de montrer que :
\[
\frac1{v_n}\to\frac1m.
\]
On écrit :
\[
\left|\frac1{v_n}-\frac1m\right|
=
\frac{\abs{v_n-m}}{\abs{m}\abs{v_n}}.
\]
Comme \(v_n\to m\neq0\), à partir d'un certain rang :
\[
\abs{v_n}\geq\frac{\abs{m}}{2}.
\]
Donc :
\[
\left|\frac1{v_n}-\frac1m\right|
\leq
\frac{2}{\abs{m}^2}\abs{v_n-m}.
\]
Le membre de droite tend vers \(0\).
\end{proof}

\begin{erreurbox}
On ne peut pas diviser par une suite qui tend vers \(0\) sans étude supplémentaire.  
La forme :
\[
\frac{0}{0}
\]
est une forme indéterminée.
\end{erreurbox}

\subsection{Composition par une fonction continue}

\begin{theorembox}
Si :
\[
u_n\to\ell
\]
et si \(f\) est continue en \(\ell\), alors :
\[
f(u_n)\to f(\ell).
\]
\end{theorembox}

\begin{proof}
C'est la traduction séquentielle de la continuité.  
Soit \(\varepsilon>0\).  
Comme \(f\) est continue en \(\ell\), il existe \(\eta>0\) tel que :
\[
\abs{x-\ell}\leq\eta \implique \abs{f(x)-f(\ell)}\leq\varepsilon.
\]
Comme \(u_n\to\ell\), il existe \(N\in\N\) tel que :
\[
n\geq N \implique \abs{u_n-\ell}\leq\eta.
\]
Donc :
\[
n\geq N \implique \abs{f(u_n)-f(\ell)}\leq\varepsilon.
\]
Ainsi :
\[
f(u_n)\to f(\ell).
\]
\end{proof}

%=========================================================
\section{Passage à la limite dans les inégalités}

\subsection{Inégalités larges}

\begin{theorembox}[Passage à la limite]
Soient \((u_n)\) et \((v_n)\) deux suites réelles convergentes telles que :
\[
\forall n\in\N,\qquad u_n\leq v_n.
\]
Si :
\[
u_n\to\ell
\quad\text{et}\quad
v_n\to m,
\]
alors :
\[
\ell\leq m.
\]
\end{theorembox}

\begin{proof}
Supposons par l'absurde que :
\[
\ell>m.
\]
Posons :
\[
\varepsilon=\frac{\ell-m}{3}>0.
\]
À partir d'un certain rang :
\[
u_n\geq \ell-\varepsilon
\]
et :
\[
v_n\leq m+\varepsilon.
\]
Or :
\[
\ell-\varepsilon
=
\ell-\frac{\ell-m}{3}
=
\frac{2\ell+m}{3},
\]
et :
\[
m+\varepsilon
=
m+\frac{\ell-m}{3}
=
\frac{\ell+2m}{3}.
\]
Comme \(\ell>m\), on a :
\[
\frac{2\ell+m}{3}>\frac{\ell+2m}{3}.
\]
Donc, pour \(n\) assez grand :
\[
u_n>v_n,
\]
ce qui contredit :
\[
u_n\leq v_n.
\]
Ainsi :
\[
\ell\leq m.
\]
\end{proof}

\begin{erreurbox}
Si :
\[
u_n<v_n
\]
pour tout \(n\), on ne peut conclure que :
\[
\lim u_n\leq \lim v_n,
\]
pas forcément une inégalité stricte.

Exemple :
\[
\frac1n>0
\]
pour tout \(n\geq1\), mais :
\[
\lim_{n\to+\infty}\frac1n=0.
\]
\end{erreurbox}

\subsection{Théorème des gendarmes}

\begin{theorembox}[Théorème des gendarmes]
Soient \((u_n)\), \((v_n)\), \((w_n)\) trois suites réelles.  
On suppose qu'à partir d'un certain rang :
\[
u_n\leq v_n\leq w_n.
\]
Si :
\[
u_n\to \ell
\quad\text{et}\quad
w_n\to \ell,
\]
alors :
\[
v_n\to \ell.
\]
\end{theorembox}

\begin{proof}
Soit \(\varepsilon>0\).  
Comme \(u_n\to\ell\), il existe \(N_1\) tel que :
\[
n\geq N_1\implique u_n\geq \ell-\varepsilon.
\]
Comme \(w_n\to\ell\), il existe \(N_2\) tel que :
\[
n\geq N_2\implique w_n\leq \ell+\varepsilon.
\]
À partir d'un rang \(N_0\), on a :
\[
u_n\leq v_n\leq w_n.
\]
Donc, pour :
\[
n\geq \max(N_0,N_1,N_2),
\]
on obtient :
\[
\ell-\varepsilon\leq u_n\leq v_n\leq w_n\leq \ell+\varepsilon.
\]
Ainsi :
\[
\abs{v_n-\ell}\leq\varepsilon.
\]
Donc :
\[
v_n\to\ell.
\]
\end{proof}

\begin{methodbox}[Utiliser le théorème des gendarmes]
Pour montrer qu'une suite \(v_n\) tend vers \(\ell\), on peut chercher deux suites plus simples \(u_n\) et \(w_n\) telles que :
\[
u_n\leq v_n\leq w_n
\]
et :
\[
u_n\to\ell,\qquad w_n\to\ell.
\]
\end{methodbox}

\begin{exercicebox}[Gendarmes]
Montrer que :
\[
\frac{\sin n}{n}\to0.
\]
\end{exercicebox}

\begin{correctionbox}
Pour tout \(n\geq1\),
\[
-1\leq\sin n\leq1.
\]
En divisant par \(n>0\), on obtient :
\[
-\frac1n\leq\frac{\sin n}{n}\leq\frac1n.
\]
Or :
\[
-\frac1n\to0
\quad\text{et}\quad
\frac1n\to0.
\]
Par le théorème des gendarmes :
\[
\frac{\sin n}{n}\to0.
\]
\end{correctionbox}

%=========================================================
\section{Suites monotones}

\subsection{Définitions}

\begin{definitionbox}[Suite croissante]
Une suite réelle \((u_n)\) est croissante si :
\[
\forall n\in\N,\qquad u_n\leq u_{n+1}.
\]
\end{definitionbox}

\begin{definitionbox}[Suite décroissante]
Une suite réelle \((u_n)\) est décroissante si :
\[
\forall n\in\N,\qquad u_{n+1}\leq u_n.
\]
\end{definitionbox}

\begin{definitionbox}[Suite monotone]
Une suite est monotone si elle est croissante ou décroissante.
\end{definitionbox}

\subsection{Méthodes d'étude de la monotonie}

\begin{methodbox}[Étudier la monotonie]
Pour étudier la monotonie de \((u_n)\), on peut :
\begin{enumerate}
    \item étudier le signe de :
    \[
    u_{n+1}-u_n ;
    \]
    \item si les termes sont strictement positifs, comparer le quotient :
    \[
    \frac{u_{n+1}}{u_n}
    \]
    à \(1\) ;
    \item si \(u_n=f(n)\), étudier la monotonie de la fonction \(f\).
\end{enumerate}
\end{methodbox}

\begin{exercicebox}[Monotonie]
Étudier la monotonie de :
\[
u_n=\frac{n}{n+1}.
\]
\end{exercicebox}

\begin{correctionbox}
On calcule :
\[
u_{n+1}-u_n
=
\frac{n+1}{n+2}-\frac{n}{n+1}.
\]
On met au même dénominateur :
\[
u_{n+1}-u_n
=
\frac{(n+1)^2-n(n+2)}{(n+1)(n+2)}.
\]
Le numérateur vaut :
\[
(n+1)^2-n(n+2)
=
n^2+2n+1-n^2-2n
=
1.
\]
Donc :
\[
u_{n+1}-u_n=\frac1{(n+1)(n+2)}>0.
\]
La suite est strictement croissante.
\end{correctionbox}

\subsection{Théorème de la limite monotone}

\begin{theorembox}[Théorème de convergence monotone]
Toute suite réelle croissante et majorée converge.  
Toute suite réelle décroissante et minorée converge.
\end{theorembox}

\begin{proof}
Soit \((u_n)\) une suite croissante et majorée.  
L'ensemble :
\[
A=\{u_n:\ n\in\N\}
\]
est non vide et majoré.  
Il admet donc une borne supérieure :
\[
\ell=\sup A.
\]
Montrons que :
\[
u_n\to\ell.
\]

Soit \(\varepsilon>0\).  
Comme \(\ell=\sup A\), il existe un élément de \(A\), donc un entier \(N\), tel que :
\[
\ell-\varepsilon<u_N\leq \ell.
\]
Comme la suite est croissante, pour tout \(n\geq N\),
\[
u_N\leq u_n\leq \ell.
\]
Donc :
\[
\ell-\varepsilon<u_n\leq \ell.
\]
Ainsi :
\[
\abs{u_n-\ell}\leq\varepsilon.
\]
Donc :
\[
u_n\to\ell.
\]

Le cas décroissant minoré se démontre de manière analogue, ou en appliquant le résultat précédent à \((-u_n)\).
\end{proof}

\begin{theorembox}[Version divergente]
Une suite croissante non majorée tend vers \(+\infty\).  
Une suite décroissante non minorée tend vers \(-\infty\).
\end{theorembox}

\begin{proof}
Soit \((u_n)\) croissante et non majorée.  
Soit \(A\in\R\).  
Comme la suite n'est pas majorée, il existe \(N\in\N\) tel que :
\[
u_N>A.
\]
Comme \((u_n)\) est croissante, pour tout \(n\geq N\),
\[
u_n\geq u_N>A.
\]
Donc :
\[
u_n\to+\infty.
\]
Le cas décroissant se traite de la même manière.
\end{proof}

\begin{retenirbox}
Le théorème de convergence monotone est une conséquence directe de la propriété de la borne supérieure de \(\R\).
\end{retenirbox}

%=========================================================
\section{Suites adjacentes}

\subsection{Définition}

\begin{definitionbox}[Suites adjacentes]
Deux suites réelles \((u_n)\) et \((v_n)\) sont adjacentes si :
\begin{enumerate}
    \item \((u_n)\) est croissante ;
    \item \((v_n)\) est décroissante ;
    \item pour tout \(n\), on a :
    \[
    u_n\leq v_n ;
    \]
    \item
    \[
    v_n-u_n\to0.
    \]
\end{enumerate}
\end{definitionbox}

\begin{theorembox}[Théorème des suites adjacentes]
Deux suites adjacentes convergent et ont la même limite.
\end{theorembox}

\begin{proof}
Comme :
\[
u_n\leq v_n
\]
pour tout \(n\), et comme \((v_n)\) est décroissante, on a :
\[
v_n\leq v_0.
\]
Donc :
\[
u_n\leq v_n\leq v_0.
\]
La suite \((u_n)\) est croissante et majorée, donc elle converge vers une limite \(\ell\).

De même, comme :
\[
u_0\leq u_n\leq v_n,
\]
la suite \((v_n)\) est minorée par \(u_0\).  
Elle est décroissante et minorée, donc elle converge vers une limite \(m\).

Comme :
\[
v_n-u_n\to0,
\]
et comme :
\[
v_n-u_n\to m-\ell,
\]
par opérations sur les limites, on obtient :
\[
m-\ell=0.
\]
Donc :
\[
m=\ell.
\]
Les deux suites convergent vers la même limite.
\end{proof}

\begin{methodbox}[Utiliser les suites adjacentes]
Pour montrer que deux suites convergent vers une même limite :
\begin{enumerate}
    \item montrer que l'une est croissante ;
    \item montrer que l'autre est décroissante ;
    \item montrer que la première reste inférieure à la seconde ;
    \item montrer que leur différence tend vers \(0\).
\end{enumerate}
\end{methodbox}

\begin{exercicebox}[Suites adjacentes]
Soient :
\[
u_n=\sum_{k=0}^n\frac1{k!},
\qquad
v_n=u_n+\frac1{n!}.
\]
Montrer que \((u_n)\) et \((v_n)\) sont adjacentes à partir d'un certain rang.
\end{exercicebox}

\begin{correctionbox}
On a :
\[
u_{n+1}-u_n=\frac1{(n+1)!}>0.
\]
Donc \((u_n)\) est croissante.

On calcule :
\[
v_{n+1}-v_n
=
\left(u_{n+1}+\frac1{(n+1)!}\right)
-
\left(u_n+\frac1{n!}\right).
\]
Or :
\[
u_{n+1}-u_n=\frac1{(n+1)!}.
\]
Donc :
\[
v_{n+1}-v_n
=
\frac1{(n+1)!}+\frac1{(n+1)!}-\frac1{n!}.
\]
Comme :
\[
\frac1{n!}=\frac{n+1}{(n+1)!},
\]
on obtient :
\[
v_{n+1}-v_n
=
\frac{2-(n+1)}{(n+1)!}
=
\frac{1-n}{(n+1)!}.
\]
Donc pour \(n\geq1\),
\[
v_{n+1}-v_n\leq0.
\]
La suite \((v_n)\) est décroissante à partir du rang \(1\).

On a évidemment :
\[
u_n\leq v_n.
\]
Enfin :
\[
v_n-u_n=\frac1{n!}\to0.
\]
Donc les deux suites sont adjacentes à partir du rang \(1\).
\end{correctionbox}

%=========================================================
\section{Suites complexes}

\subsection{Convergence d'une suite complexe}

\begin{definitionbox}[Suite complexe convergente]
Soit \((z_n)\) une suite complexe et soit \(\ell\in\C\).  
On dit que :
\[
z_n\to\ell
\]
si :
\[
\forall\varepsilon>0,\quad \exists N\in\N,\quad \forall n\geq N,\quad \abs{z_n-\ell}\leq\varepsilon.
\]
\end{definitionbox}

La définition est la même que pour les suites réelles, mais \(\abs{\cdot}\) désigne ici le module complexe.

\subsection{Parties réelle et imaginaire}

\begin{theorembox}
Soit :
\[
z_n=a_n+ib_n,
\qquad
\ell=\alpha+i\beta.
\]
Alors :
\[
z_n\to\ell
\]
si et seulement si :
\[
a_n\to\alpha
\quad\text{et}\quad
b_n\to\beta.
\]
\end{theorembox}

\begin{proof}
Supposons :
\[
z_n\to\ell.
\]
Alors :
\[
\abs{a_n-\alpha}\leq \abs{z_n-\ell},
\]
et :
\[
\abs{b_n-\beta}\leq \abs{z_n-\ell}.
\]
Donc :
\[
a_n\to\alpha
\quad\text{et}\quad
b_n\to\beta.
\]

Réciproquement, supposons :
\[
a_n\to\alpha
\quad\text{et}\quad
b_n\to\beta.
\]
Alors :
\[
\abs{z_n-\ell}
=
\sqrt{(a_n-\alpha)^2+(b_n-\beta)^2}.
\]
On a :
\[
\abs{z_n-\ell}
\leq
\abs{a_n-\alpha}+\abs{b_n-\beta}.
\]
Le membre de droite tend vers \(0\).  
Donc :
\[
z_n\to\ell.
\]
\end{proof}

\begin{retenirbox}
Une suite complexe converge si et seulement si ses parties réelle et imaginaire convergent.
\end{retenirbox}

%=========================================================
\section{Suites extraites}

\subsection{Définition}

\begin{definitionbox}[Suite extraite]
Soit \((u_n)\) une suite.  
Une suite extraite de \((u_n)\) est une suite de la forme :
\[
(u_{\varphi(n)})_{n\in\N},
\]
où :
\[
\varphi:\N\to\N
\]
est strictement croissante.
\end{definitionbox}

Exemples :
\[
(u_{2n})_{n\in\N}
\]
est la suite des termes d'indice pair.

\[
(u_{2n+1})_{n\in\N}
\]
est la suite des termes d'indice impair.

\begin{theorembox}
Si \(\varphi:\N\to\N\) est strictement croissante, alors :
\[
\forall n\in\N,\qquad \varphi(n)\geq n.
\]
\end{theorembox}

\begin{proof}
On raisonne par récurrence.  
Pour \(n=0\), on a :
\[
\varphi(0)\geq0.
\]
Supposons :
\[
\varphi(n)\geq n.
\]
Comme \(\varphi\) est strictement croissante :
\[
\varphi(n+1)>\varphi(n).
\]
Comme \(\varphi(n)\) et \(\varphi(n+1)\) sont des entiers :
\[
\varphi(n+1)\geq\varphi(n)+1.
\]
Donc :
\[
\varphi(n+1)\geq n+1.
\]
\end{proof}

\subsection{Suites extraites d'une suite convergente}

\begin{theorembox}
Si :
\[
u_n\to\ell,
\]
alors toute suite extraite :
\[
u_{\varphi(n)}
\]
converge aussi vers \(\ell\).
\end{theorembox}

\begin{proof}
Soit \(\varepsilon>0\).  
Comme \(u_n\to\ell\), il existe \(N\in\N\) tel que :
\[
k\geq N\implique \abs{u_k-\ell}\leq\varepsilon.
\]
Comme \(\varphi(n)\geq n\), pour tout \(n\geq N\),
\[
\varphi(n)\geq n\geq N.
\]
Donc :
\[
\abs{u_{\varphi(n)}-\ell}\leq\varepsilon.
\]
Ainsi :
\[
u_{\varphi(n)}\to\ell.
\]
\end{proof}

\begin{methodbox}[Montrer qu'une suite diverge avec deux sous-suites]
Pour montrer qu'une suite \((u_n)\) ne converge pas, il suffit de trouver deux suites extraites convergeant vers deux limites différentes.
\end{methodbox}

\begin{exercicebox}[Divergence par extraction]
Montrer que la suite :
\[
u_n=(-1)^n
\]
ne converge pas.
\end{exercicebox}

\begin{correctionbox}
Considérons les deux suites extraites :
\[
u_{2n}=1
\]
et :
\[
u_{2n+1}=-1.
\]
La première converge vers \(1\), la seconde converge vers \(-1\).  
Si \((u_n)\) convergeait, toutes ses suites extraites convergeraient vers la même limite.  
Ce n'est pas le cas.  
Donc \((u_n)\) diverge.
\end{correctionbox}

\subsection{Recouvrement par sous-suites}

\begin{theorembox}
Si :
\[
u_{2n}\to\ell
\quad\text{et}\quad
u_{2n+1}\to\ell,
\]
alors :
\[
u_n\to\ell.
\]
\end{theorembox}

\begin{proof}
Soit \(\varepsilon>0\).  
Il existe \(N_1\) tel que :
\[
n\geq N_1\implique \abs{u_{2n}-\ell}\leq\varepsilon.
\]
Il existe \(N_2\) tel que :
\[
n\geq N_2\implique \abs{u_{2n+1}-\ell}\leq\varepsilon.
\]
Posons :
\[
N=2\max(N_1,N_2)+1.
\]
Pour tout \(k\geq N\), soit \(k\) est pair, soit \(k\) est impair.  
Dans les deux cas, l'indice correspondant est au moins \(\max(N_1,N_2)\).  
Donc :
\[
\abs{u_k-\ell}\leq\varepsilon.
\]
Ainsi :
\[
u_n\to\ell.
\]
\end{proof}

%=========================================================
\section{Théorème de Bolzano-Weierstrass}

\subsection{Énoncé}

\begin{theorembox}[Bolzano-Weierstrass]
De toute suite réelle bornée, on peut extraire une sous-suite convergente.
\end{theorembox}

\begin{proof}
Soit \((u_n)\) une suite réelle bornée.  
Il existe donc un segment :
\[
[a,b]
\]
tel que :
\[
\forall n\in\N,\qquad u_n\in[a,b].
\]

On construit par dichotomie une suite de segments emboîtés contenant chacun une infinité de termes de la suite.

On pose :
\[
I_0=[a,b].
\]
On coupe \(I_0\) en deux segments de même longueur.  
Au moins l'un des deux contient une infinité de termes de la suite.  
On l'appelle \(I_1\).

On recommence : si \(I_k\) contient une infinité de termes de la suite, on le coupe en deux.  
L'un des deux sous-segments contient une infinité de termes de la suite.  
On l'appelle \(I_{k+1}\).

On obtient ainsi une suite de segments emboîtés :
\[
I_0\supset I_1\supset I_2\supset \cdots
\]
dont la longueur vaut :
\[
\frac{b-a}{2^k}.
\]
Elle tend vers \(0\).

On choisit ensuite une sous-suite \((u_{\varphi(k)})\) telle que :
\[
u_{\varphi(k)}\in I_k,
\]
avec :
\[
\varphi(k+1)>\varphi(k).
\]
C'est possible parce que chaque \(I_k\) contient une infinité de termes de la suite.

Par le théorème des segments emboîtés, l'intersection des \(I_k\) est réduite à un seul réel \(\ell\).  
Comme :
\[
u_{\varphi(k)}\in I_k
\]
et que la longueur de \(I_k\) tend vers \(0\), on obtient :
\[
u_{\varphi(k)}\to\ell.
\]
Donc la suite admet une sous-suite convergente.
\end{proof}

\begin{retenirbox}
Bolzano-Weierstrass est un théorème d'extraction :  
\[
\text{bornée}
\quad\Longrightarrow\quad
\text{possède une sous-suite convergente}.
\]
Il ne dit pas que la suite entière converge.
\end{retenirbox}

\begin{erreurbox}
Une suite bornée n'est pas forcément convergente.  
Bolzano-Weierstrass affirme seulement qu'elle possède une sous-suite convergente.
\[
(-1)^n
\]
est bornée et possède deux sous-suites convergentes, mais elle ne converge pas.
\end{erreurbox}
%=========================================================
\section{Caractérisation séquentielle des bornes supérieure et inférieure}

\subsection{Motivation}

Dans le chapitre précédent, on a défini la borne supérieure d'une partie non vide et majorée de \(\R\) comme le plus petit de ses majorants.

On a également vu la caractérisation suivante :
\[
s=\sup A
\]
si et seulement si :
\[
\forall x\in A,\quad x\leq s,
\]
et :
\[
\forall \varepsilon>0,\quad \exists x\in A,\quad s-\varepsilon<x\leq s.
\]

Cette caractérisation affirme que l'on peut trouver dans \(A\) des éléments arbitrairement proches de \(s\), par valeurs inférieures.

L'idée naturelle est alors de construire une suite d'éléments de \(A\) qui converge vers \(s\).

\begin{retenirbox}
La borne supérieure n'appartient pas forcément à l'ensemble, mais elle peut toujours être approchée par une suite d'éléments de l'ensemble.
\end{retenirbox}

%---------------------------------------------------------
\subsection{Caractérisation séquentielle de la borne supérieure}

\begin{theorembox}[Caractérisation séquentielle de la borne supérieure]
Soit \(A\) une partie non vide et majorée de \(\R\), et soit \(s\in\R\).  
Alors :
\[
s=\sup A
\]
si et seulement si les deux conditions suivantes sont vérifiées :
\begin{enumerate}
    \item \(s\) est un majorant de \(A\) :
    \[
    \forall x\in A,\quad x\leq s ;
    \]
    \item il existe une suite \((a_n)_{n\geq1}\) d'éléments de \(A\) telle que :
    \[
    a_n\longrightarrow s.
    \]
\end{enumerate}
\end{theorembox}

\begin{proof}
Supposons d'abord :
\[
s=\sup A.
\]
Alors \(s\) est un majorant de \(A\).

De plus, par la caractérisation de la borne supérieure, pour tout \(n\geq1\), en prenant :
\[
\varepsilon=\frac1n,
\]
il existe un élément \(a_n\in A\) tel que :
\[
s-\frac1n<a_n\leq s.
\]
On obtient donc une suite \((a_n)\) d'éléments de \(A\) vérifiant :
\[
s-\frac1n<a_n\leq s.
\]
Ainsi :
\[
0\leq s-a_n<\frac1n.
\]
Comme :
\[
\frac1n\longrightarrow0,
\]
le théorème des gendarmes donne :
\[
s-a_n\longrightarrow0.
\]
Donc :
\[
a_n\longrightarrow s.
\]

Réciproquement, supposons que \(s\) soit un majorant de \(A\) et qu'il existe une suite \((a_n)\) d'éléments de \(A\) telle que :
\[
a_n\to s.
\]
Montrons que \(s\) est le plus petit majorant de \(A\).

Soit \(M<s\).  
On veut montrer que \(M\) n'est pas un majorant de \(A\).  
Posons :
\[
\varepsilon=s-M>0.
\]
Comme :
\[
a_n\to s,
\]
il existe \(N\in\N\) tel que, pour tout \(n\geq N\),
\[
|a_n-s|<\varepsilon.
\]
En particulier :
\[
s-\varepsilon<a_N.
\]
Or :
\[
s-\varepsilon=s-(s-M)=M.
\]
Donc :
\[
M<a_N.
\]
Comme :
\[
a_N\in A,
\]
le réel \(M\) ne majore pas \(A\).

Ainsi, tout réel strictement inférieur à \(s\) n'est pas majorant de \(A\).  
Comme \(s\) est un majorant, \(s\) est le plus petit majorant de \(A\).  
Donc :
\[
s=\sup A.
\]
\end{proof}

\begin{retenirbox}
Pour montrer que \(s=\sup A\), on peut montrer :
\[
\forall x\in A,\quad x\leq s
\]
puis construire une suite \((a_n)\) d'éléments de \(A\) telle que :
\[
a_n\to s.
\]
\end{retenirbox}

%---------------------------------------------------------
\subsection{Caractérisation séquentielle de la borne inférieure}

\begin{theorembox}[Caractérisation séquentielle de la borne inférieure]
Soit \(A\) une partie non vide et minorée de \(\R\), et soit \(i\in\R\).  
Alors :
\[
i=\inf A
\]
si et seulement si les deux conditions suivantes sont vérifiées :
\begin{enumerate}
    \item \(i\) est un minorant de \(A\) :
    \[
    \forall x\in A,\quad i\leq x ;
    \]
    \item il existe une suite \((a_n)_{n\geq1}\) d'éléments de \(A\) telle que :
    \[
    a_n\longrightarrow i.
    \]
\end{enumerate}
\end{theorembox}

\begin{proof}
On raisonne de manière analogue au cas de la borne supérieure.

Supposons :
\[
i=\inf A.
\]
Alors \(i\) est un minorant de \(A\).  
De plus, pour tout \(n\geq1\), par caractérisation de la borne inférieure, il existe \(a_n\in A\) tel que :
\[
i\leq a_n<i+\frac1n.
\]
On obtient :
\[
0\leq a_n-i<\frac1n.
\]
Par le théorème des gendarmes :
\[
a_n-i\to0.
\]
Donc :
\[
a_n\to i.
\]

Réciproquement, supposons que \(i\) soit un minorant de \(A\) et qu'il existe une suite \((a_n)\) d'éléments de \(A\) telle que :
\[
a_n\to i.
\]
Montrons que \(i\) est le plus grand minorant de \(A\).

Soit \(m>i\).  
On veut montrer que \(m\) n'est pas un minorant de \(A\).  
Posons :
\[
\varepsilon=m-i>0.
\]
Comme :
\[
a_n\to i,
\]
il existe \(N\in\N\) tel que :
\[
|a_N-i|<\varepsilon.
\]
Donc :
\[
a_N<i+\varepsilon=m.
\]
Comme :
\[
a_N\in A,
\]
le réel \(m\) ne minore pas \(A\).

Ainsi, tout réel strictement supérieur à \(i\) n'est pas minorant de \(A\).  
Comme \(i\) est un minorant, \(i\) est le plus grand minorant de \(A\).  
Donc :
\[
i=\inf A.
\]
\end{proof}

%---------------------------------------------------------
\subsection{Exemple fondamental : intervalle ouvert}

\begin{exercicebox}[Borne supérieure et suite approchante]
Soit :
\[
A=]0,1[.
\]
Construire une suite d'éléments de \(A\) qui converge vers \(\sup A\).
\end{exercicebox}

\begin{correctionbox}
On sait que :
\[
\sup A=1.
\]
Pour construire une suite d'éléments de \(A\) qui converge vers \(1\), on peut prendre :
\[
a_n=1-\frac1{n+1}.
\]
Alors, pour tout \(n\geq1\),
\[
0<a_n<1,
\]
donc :
\[
a_n\in A.
\]
De plus :
\[
a_n=1-\frac1{n+1}\longrightarrow1.
\]
Ainsi, \((a_n)\) est bien une suite d'éléments de \(A\) qui converge vers :
\[
\sup A=1.
\]
\end{correctionbox}

\begin{exercicebox}[Borne inférieure et suite approchante]
Soit :
\[
A=]0,1[.
\]
Construire une suite d'éléments de \(A\) qui converge vers \(\inf A\).
\end{exercicebox}

\begin{correctionbox}
On sait que :
\[
\inf A=0.
\]
On peut prendre :
\[
a_n=\frac1{n+1}.
\]
Alors, pour tout \(n\geq1\),
\[
0<a_n<1,
\]
donc :
\[
a_n\in A.
\]
De plus :
\[
a_n\longrightarrow0.
\]
Ainsi, \((a_n)\) est une suite d'éléments de \(A\) qui converge vers :
\[
\inf A=0.
\]
\end{correctionbox}

%---------------------------------------------------------
\subsection{Exemple avec un ensemble discret}

\begin{exercicebox}[Ensemble de la forme \(\{1-\frac1n\}\)]
Soit :
\[
A=\left\{1-\frac1n:\ n\in\N^*\right\}.
\]
Déterminer \(\sup A\), \(\inf A\), et construire des suites d'éléments de \(A\) qui convergent vers ces bornes lorsque c'est possible.
\end{exercicebox}

\begin{correctionbox}
Les premiers éléments de \(A\) sont :
\[
0,\quad \frac12,\quad \frac23,\quad \frac34,\dots
\]
On a :
\[
1-\frac1n<1
\]
pour tout \(n\in\N^*\).  
Donc \(1\) est un majorant de \(A\).

La suite :
\[
a_n=1-\frac1n
\]
est une suite d'éléments de \(A\) et :
\[
a_n\to1.
\]
Donc, par caractérisation séquentielle :
\[
\sup A=1.
\]

Pour la borne inférieure, on remarque que :
\[
0\in A
\]
car :
\[
1-\frac11=0.
\]
De plus, tous les éléments de \(A\) sont positifs ou nuls :
\[
1-\frac1n\geq0.
\]
Donc :
\[
\min A=0.
\]
Ainsi :
\[
\inf A=0.
\]

Une suite d'éléments de \(A\) qui converge vers \(\inf A=0\) peut être constante :
\[
b_n=0.
\]
En effet :
\[
b_n\in A
\]
pour tout \(n\), et :
\[
b_n\to0.
\]
\end{correctionbox}

%---------------------------------------------------------
\subsection{Méthode générale}

\begin{methodbox}[Construire une suite approchant une borne supérieure]
Soit \(A\subset\R\) non vide et majorée.  
Pour construire une suite \((a_n)\) d'éléments de \(A\) qui converge vers \(\sup A\), on utilise :
\[
\forall n\geq1,\quad \exists a_n\in A,\quad
\sup A-\frac1n<a_n\leq \sup A.
\]
Alors :
\[
a_n\to\sup A.
\]
\end{methodbox}

\begin{methodbox}[Construire une suite approchant une borne inférieure]
Soit \(A\subset\R\) non vide et minorée.  
Pour construire une suite \((a_n)\) d'éléments de \(A\) qui converge vers \(\inf A\), on utilise :
\[
\forall n\geq1,\quad \exists a_n\in A,\quad
\inf A\leq a_n<\inf A+\frac1n.
\]
Alors :
\[
a_n\to\inf A.
\]
\end{methodbox}

\begin{erreurbox}
La suite approchante n'est pas forcément unique.  
Il peut exister de très nombreuses suites d'éléments de \(A\) qui convergent vers \(\sup A\) ou \(\inf A\).
\end{erreurbox}

\begin{erreurbox}
Le fait qu'il existe une suite \((a_n)\) d'éléments de \(A\) telle que \(a_n\to s\) ne suffit pas à prouver que :
\[
s=\sup A.
\]
Il faut aussi prouver que \(s\) est un majorant de \(A\).
\]
Par exemple, pour :
\[
A=[0,2],
\]
la suite constante :
\[
a_n=1
\]
converge vers \(1\), mais :
\[
1\neq\sup A.
\]
\end{erreurbox}

%---------------------------------------------------------
\subsection{Application : suites monotones}

\begin{theorembox}[Limite d'une suite monotone avec borne supérieure]
Soit \((u_n)\) une suite réelle croissante et majorée.  
Alors :
\[
u_n\to \sup\{u_n:\ n\in\N\}.
\]
\end{theorembox}

\begin{proof}
Posons :
\[
A=\{u_n:\ n\in\N\}.
\]
Comme \((u_n)\) est majorée, \(A\) est une partie non vide et majorée de \(\R\).  
Elle admet donc une borne supérieure :
\[
\ell=\sup A.
\]
Comme \(\ell=\sup A\), pour tout \(\varepsilon>0\), il existe un entier \(N\) tel que :
\[
\ell-\varepsilon<u_N\leq \ell.
\]
Comme la suite est croissante, pour tout \(n\geq N\),
\[
u_N\leq u_n\leq \ell.
\]
Donc :
\[
\ell-\varepsilon<u_n\leq \ell.
\]
Ainsi :
\[
|u_n-\ell|\leq\varepsilon.
\]
Donc :
\[
u_n\to\ell.
\]
\end{proof}

\begin{theorembox}[Limite d'une suite monotone avec borne inférieure]
Soit \((u_n)\) une suite réelle décroissante et minorée.  
Alors :
\[
u_n\to \inf\{u_n:\ n\in\N\}.
\]
\end{theorembox}

\begin{proof}
La preuve est analogue.  
On pose :
\[
A=\{u_n:\ n\in\N\}.
\]
Comme \((u_n)\) est décroissante et minorée, \(A\) est non vide et minorée.  
On note :
\[
\ell=\inf A.
\]
Pour tout \(\varepsilon>0\), il existe \(N\) tel que :
\[
\ell\leq u_N<\ell+\varepsilon.
\]
Comme la suite est décroissante, pour tout \(n\geq N\),
\[
\ell\leq u_n\leq u_N<\ell+\varepsilon.
\]
Donc :
\[
|u_n-\ell|\leq\varepsilon.
\]
Ainsi :
\[
u_n\to\ell.
\]
\end{proof}

%=========================================================
\section{Suites récurrentes}

\subsection{Suites de la forme \(u_{n+1}=f(u_n)\)}

\begin{definitionbox}[Suite récurrente d'ordre 1]
Une suite récurrente d'ordre \(1\) est une suite définie par :
\[
u_0\in I,
\qquad
u_{n+1}=f(u_n),
\]
où \(f:I\to I\).
\end{definitionbox}

\begin{retenirbox}
Pour étudier une suite récurrente :
\[
u_{n+1}=f(u_n),
\]
il faut souvent étudier la fonction \(f\).
\end{retenirbox}

\subsection{Stabilité de l'intervalle}

\begin{methodbox}[Montrer que la suite est bien définie]
Pour montrer que la suite est bien définie dans un intervalle \(I\), on montre :
\[
u_0\in I
\]
et :
\[
f(I)\subset I.
\]
Alors, par récurrence :
\[
\forall n,\quad u_n\in I.
\]
\end{methodbox}

\subsection{Limite éventuelle}

\begin{theorembox}[Limite d'une suite récurrente]
Soit \(f:I\to I\) continue.  
Si la suite \((u_n)\) définie par :
\[
u_{n+1}=f(u_n)
\]
converge vers \(\ell\in I\), alors :
\[
f(\ell)=\ell.
\]
\end{theorembox}

\begin{proof}
Si :
\[
u_n\to\ell,
\]
alors :
\[
u_{n+1}\to\ell.
\]
Comme \(f\) est continue :
\[
f(u_n)\to f(\ell).
\]
Or :
\[
u_{n+1}=f(u_n).
\]
Donc, par unicité de la limite :
\[
\ell=f(\ell).
\]
\end{proof}

\begin{retenirbox}
La limite éventuelle d'une suite récurrente est un point fixe de \(f\).
\[
\ell=f(\ell).
\]
\end{retenirbox}

\begin{erreurbox}
Résoudre :
\[
f(\ell)=\ell
\]
ne prouve pas que la suite converge.  
Cela donne seulement les limites possibles.
\end{erreurbox}

\subsection{Étude par monotonie}

\begin{methodbox}[Suite récurrente monotone]
Pour étudier la monotonie de :
\[
u_{n+1}=f(u_n),
\]
on peut étudier le signe de :
\[
u_{n+1}-u_n=f(u_n)-u_n.
\]
On introduit souvent :
\[
g(x)=f(x)-x.
\]
\end{methodbox}

\begin{exercicebox}[Suite récurrente classique]
Soit :
\[
u_0>0,
\qquad
u_{n+1}=\frac12\left(u_n+\frac2{u_n}\right).
\]
Montrer que, si \(u_0\geq \sqrt2\), la suite converge vers \(\sqrt2\).
\end{exercicebox}

\begin{correctionbox}
On suppose :
\[
u_0\geq\sqrt2.
\]
Montrons d'abord que :
\[
u_n\geq\sqrt2
\]
pour tout \(n\).  
Si \(u_n>0\), alors :
\[
u_{n+1}-\sqrt2
=
\frac12\left(u_n+\frac2{u_n}\right)-\sqrt2.
\]
On met au même dénominateur :
\[
u_{n+1}-\sqrt2
=
\frac{u_n^2+2-2\sqrt2u_n}{2u_n}
=
\frac{(u_n-\sqrt2)^2}{2u_n}\geq0.
\]
Donc :
\[
u_{n+1}\geq\sqrt2.
\]
Par récurrence :
\[
u_n\geq\sqrt2.
\]

Étudions la monotonie :
\[
u_{n+1}-u_n
=
\frac12\left(u_n+\frac2{u_n}\right)-u_n
=
\frac{2-u_n^2}{2u_n}.
\]
Comme :
\[
u_n\geq\sqrt2,
\]
on a :
\[
2-u_n^2\leq0.
\]
Donc :
\[
u_{n+1}-u_n\leq0.
\]
La suite est décroissante et minorée par \(\sqrt2\).  
Elle converge donc vers une limite \(\ell\geq\sqrt2\).

La fonction :
\[
f(x)=\frac12\left(x+\frac2x\right)
\]
est continue sur \(\R_+^*\).  
En passant à la limite dans :
\[
u_{n+1}=f(u_n),
\]
on obtient :
\[
\ell=\frac12\left(\ell+\frac2\ell\right).
\]
Donc :
\[
2\ell=\ell+\frac2\ell.
\]
Ainsi :
\[
\ell=\frac2\ell,
\]
donc :
\[
\ell^2=2.
\]
Comme :
\[
\ell\geq0,
\]
on obtient :
\[
\ell=\sqrt2.
\]
\end{correctionbox}

%=========================================================
\section{Comparaison de suites}

\subsection{Négligeabilité}

\begin{definitionbox}[Négligeabilité]
Soient \((u_n)\) et \((v_n)\) deux suites réelles ou complexes, avec \(v_n\neq0\) à partir d'un certain rang.  
On dit que \(u_n\) est négligeable devant \(v_n\) si :
\[
\frac{u_n}{v_n}\to0.
\]
On note :
\[
u_n=o(v_n).
\]
\end{definitionbox}

Exemple :
\[
1=o(n),
\qquad
n=o(n^2).
\]

\subsection{Domination}

\begin{definitionbox}[Domination]
On dit que \(u_n\) est dominée par \(v_n\) si la suite :
\[
\left(\frac{u_n}{v_n}\right)
\]
est bornée à partir d'un certain rang.  
On note :
\[
u_n=O(v_n).
\]
\end{definitionbox}

Exemple :
\[
3n^2+2n+1=O(n^2).
\]

\subsection{Équivalence}

\begin{definitionbox}[Équivalence]
Soient \((u_n)\) et \((v_n)\) deux suites, avec \(v_n\neq0\) à partir d'un certain rang.  
On dit que \(u_n\) est équivalente à \(v_n\) si :
\[
\frac{u_n}{v_n}\to1.
\]
On note :
\[
u_n\sim v_n.
\]
\end{definitionbox}

\begin{retenirbox}
\[
u_n\sim v_n
\]
signifie que \(u_n\) et \(v_n\) ont le même ordre de grandeur asymptotique.
\end{retenirbox}

\begin{theorembox}
Si :
\[
u_n\sim v_n
\]
et si :
\[
v_n\to\ell
\]
avec \(\ell\in\R\), alors :
\[
u_n\to\ell.
\]
\end{theorembox}

\begin{proof}
Comme :
\[
u_n\sim v_n,
\]
on peut écrire :
\[
u_n=v_n\varepsilon_n
\]
avec :
\[
\varepsilon_n\to1.
\]
Si :
\[
v_n\to\ell,
\]
alors par produit :
\[
u_n=v_n\varepsilon_n\to \ell\cdot1=\ell.
\]
\end{proof}

\begin{erreurbox}
On ne peut pas additionner directement des équivalents :
\[
u_n\sim v_n
\quad\text{et}\quad
a_n\sim b_n
\]
n'implique pas toujours :
\[
u_n+a_n\sim v_n+b_n.
\]
Il peut y avoir compensation.
\end{erreurbox}

\subsection{Équivalents usuels de suites simples}

\begin{theorembox}[Équivalents classiques]
Lorsque \(n\to+\infty\), on a :
\[
n+1\sim n,
\]
\[
an^p+o(n^p)\sim an^p
\quad(a\neq0),
\]
\[
\frac{P(n)}{Q(n)}
\sim
\frac{a_p n^p}{b_q n^q}
=
\frac{a_p}{b_q}n^{p-q},
\]
où \(P\) et \(Q\) sont des polynômes de degrés respectifs \(p\) et \(q\), de coefficients dominants \(a_p\) et \(b_q\).
\end{theorembox}

\begin{exercicebox}[Équivalent simple]
Donner un équivalent de :
\[
u_n=\frac{3n^2+2n+1}{5n^3-n}.
\]
\end{exercicebox}

\begin{correctionbox}
Au numérateur, le terme dominant est :
\[
3n^2.
\]
Au dénominateur, le terme dominant est :
\[
5n^3.
\]
Donc :
\[
u_n\sim \frac{3n^2}{5n^3}
=
\frac{3}{5n}.
\]
\end{correctionbox}

%=========================================================
\newpage
\section{Exercices progressifs}

\subsection*{Thème 1 -- Convergence et limites}

\begin{exercicebox}[1]
Montrer avec la définition que :
\[
\frac{2n+1}{n}\to2.
\]
\end{exercicebox}

\begin{exercicebox}[2]
Montrer que :
\[
\frac{n^2+1}{n^2+n}\to1.
\]
\end{exercicebox}

\begin{exercicebox}[3]
Montrer que :
\[
\frac{(-1)^n}{n}\to0.
\]
\end{exercicebox}

\subsection*{Thème 2 -- Monotonie}

\begin{exercicebox}[4]
Étudier la monotonie et la convergence de :
\[
u_n=\frac{n}{n+2}.
\]
\end{exercicebox}

\begin{exercicebox}[5]
Soit :
\[
u_0=1,
\qquad
u_{n+1}=\frac{u_n+3}{2}.
\]
Étudier la convergence de \((u_n)\).
\end{exercicebox}

\subsection*{Thème 3 -- Suites extraites}

\begin{exercicebox}[6]
Montrer que la suite :
\[
u_n=\cos(n\pi)
\]
ne converge pas.
\end{exercicebox}

\begin{exercicebox}[7]
Soit \((u_n)\) une suite telle que :
\[
u_{2n}\to\ell
\quad\text{et}\quad
u_{2n+1}\to\ell.
\]
Montrer que :
\[
u_n\to\ell.
\]
\end{exercicebox}

\subsection*{Thème 4 -- Comparaisons}

\begin{exercicebox}[8]
Donner un équivalent de :
\[
u_n=\frac{2n^3-n+1}{7n^2+5}.
\]
\end{exercicebox}

\begin{exercicebox}[9]
Montrer que :
\[
\sqrt{n^2+n}\sim n.
\]
\end{exercicebox}

\subsection*{Thème 5 -- Bilan}

\begin{exercicebox}[10]
Soit :
\[
u_n=\sum_{k=1}^n\frac1{k^2}.
\]
Montrer que \((u_n)\) est croissante et majorée.
\end{exercicebox}

%=========================================================
\newpage
\section{Corrigés détaillés des exercices progressifs}

\begin{correctionbox}[1]
On a :
\[
\frac{2n+1}{n}=2+\frac1n.
\]
Donc :
\[
\left|\frac{2n+1}{n}-2\right|=\frac1n.
\]
Soit \(\varepsilon>0\).  
On choisit \(N\in\N^*\) tel que :
\[
N\geq\frac1\varepsilon.
\]
Alors, pour tout \(n\geq N\),
\[
\frac1n\leq\frac1N\leq\varepsilon.
\]
Donc :
\[
\frac{2n+1}{n}\to2.
\]
\end{correctionbox}

\begin{correctionbox}[2]
On écrit :
\[
\frac{n^2+1}{n^2+n}
=
\frac{1+\frac1{n^2}}{1+\frac1n}.
\]
Comme :
\[
\frac1{n^2}\to0
\quad\text{et}\quad
\frac1n\to0,
\]
on obtient par opérations sur les limites :
\[
\frac{1+\frac1{n^2}}{1+\frac1n}\to\frac{1}{1}=1.
\]
\end{correctionbox}

\begin{correctionbox}[3]
On a :
\[
\left|\frac{(-1)^n}{n}\right|=\frac1n.
\]
Comme :
\[
-\frac1n\leq\frac{(-1)^n}{n}\leq\frac1n,
\]
et que :
\[
-\frac1n\to0,\qquad \frac1n\to0,
\]
le théorème des gendarmes donne :
\[
\frac{(-1)^n}{n}\to0.
\]
\end{correctionbox}

\begin{correctionbox}[4]
On considère :
\[
u_n=\frac{n}{n+2}.
\]
On calcule :
\[
u_{n+1}-u_n
=
\frac{n+1}{n+3}-\frac{n}{n+2}.
\]
On met au même dénominateur :
\[
u_{n+1}-u_n
=
\frac{(n+1)(n+2)-n(n+3)}{(n+3)(n+2)}.
\]
Le numérateur vaut :
\[
n^2+3n+2-n^2-3n=2.
\]
Donc :
\[
u_{n+1}-u_n=\frac{2}{(n+2)(n+3)}>0.
\]
La suite est croissante.

De plus :
\[
u_n=\frac{n}{n+2}=\frac1{1+\frac2n}\to1.
\]
Donc la suite converge vers \(1\).
\end{correctionbox}

\begin{correctionbox}[5]
On cherche d'abord la limite éventuelle.  
Si \(u_n\to\ell\), alors :
\[
\ell=\frac{\ell+3}{2}.
\]
Donc :
\[
2\ell=\ell+3,
\]
et :
\[
\ell=3.
\]

Étudions :
\[
u_{n+1}-3
=
\frac{u_n+3}{2}-3
=
\frac{u_n-3}{2}.
\]
Posons :
\[
v_n=u_n-3.
\]
Alors :
\[
v_{n+1}=\frac12v_n.
\]
Donc :
\[
v_n=\left(\frac12\right)^n v_0.
\]
Or :
\[
v_0=1-3=-2.
\]
Donc :
\[
u_n-3=-2\left(\frac12\right)^n.
\]
Ainsi :
\[
u_n=3-2\left(\frac12\right)^n.
\]
Comme :
\[
\left(\frac12\right)^n\to0,
\]
on obtient :
\[
u_n\to3.
\]
\end{correctionbox}

\begin{correctionbox}[6]
On a :
\[
u_n=\cos(n\pi)=(-1)^n.
\]
Alors :
\[
u_{2n}=1
\]
et :
\[
u_{2n+1}=-1.
\]
Les deux sous-suites convergent vers deux limites différentes.  
Donc \((u_n)\) ne converge pas.
\end{correctionbox}

\begin{correctionbox}[7]
Soit \(\varepsilon>0\).  
Il existe \(N_1\) tel que :
\[
n\geq N_1\implique \abs{u_{2n}-\ell}\leq\varepsilon.
\]
Il existe \(N_2\) tel que :
\[
n\geq N_2\implique \abs{u_{2n+1}-\ell}\leq\varepsilon.
\]
Posons :
\[
N=2\max(N_1,N_2)+1.
\]
Si \(k\geq N\), alors \(k\) est soit pair, soit impair.  
Dans les deux cas, il s'écrit :
\[
k=2n
\quad\text{ou}\quad
k=2n+1
\]
avec :
\[
n\geq \max(N_1,N_2).
\]
Donc :
\[
\abs{u_k-\ell}\leq\varepsilon.
\]
Ainsi :
\[
u_n\to\ell.
\]
\end{correctionbox}

\begin{correctionbox}[8]
Au numérateur :
\[
2n^3-n+1\sim2n^3.
\]
Au dénominateur :
\[
7n^2+5\sim7n^2.
\]
Donc :
\[
u_n\sim\frac{2n^3}{7n^2}
=
\frac27 n.
\]
\end{correctionbox}

\begin{correctionbox}[9]
On écrit :
\[
\frac{\sqrt{n^2+n}}{n}
=
\sqrt{1+\frac1n}.
\]
Comme :
\[
1+\frac1n\to1
\]
et que la fonction racine carrée est continue en \(1\), on obtient :
\[
\sqrt{1+\frac1n}\to1.
\]
Donc :
\[
\sqrt{n^2+n}\sim n.
\]
\end{correctionbox}

\begin{correctionbox}[10]
On a :
\[
u_{n+1}-u_n=\frac1{(n+1)^2}>0.
\]
Donc la suite \((u_n)\) est croissante.

Montrons qu'elle est majorée.  
Pour tout \(k\geq2\),
\[
\frac1{k^2}\leq \frac1{k(k-1)}.
\]
Or :
\[
\frac1{k(k-1)}=\frac1{k-1}-\frac1k.
\]
Donc, pour \(n\geq2\),
\[
\sum_{k=2}^n\frac1{k^2}
\leq
\sum_{k=2}^n\left(\frac1{k-1}-\frac1k\right).
\]
La somme télescope :
\[
\sum_{k=2}^n\left(\frac1{k-1}-\frac1k\right)
=
1-\frac1n.
\]
Ainsi :
\[
u_n
=
1+\sum_{k=2}^n\frac1{k^2}
\leq
1+1-\frac1n
<2.
\]
Donc \((u_n)\) est majorée par \(2\).  
Elle est croissante et majorée, donc convergente.
\end{correctionbox}

%=========================================================
\newpage
\section{Grand devoir bilan final}

\begin{center}
\Large\textbf{Devoir bilan -- Suites numériques}
\end{center}

\subsection*{Exercice 1 -- Questions de cours}

\begin{enumerate}
    \item Donner la définition de :
    \[
    u_n\to\ell.
    \]
    \item Démontrer l'unicité de la limite.
    \item Démontrer que toute suite convergente est bornée.
    \item Énoncer et démontrer le théorème des gendarmes.
    \item Énoncer le théorème de convergence monotone.
\end{enumerate}

\subsection*{Exercice 2 -- Limites}

Déterminer les limites suivantes :
\[
u_n=\frac{3n^2+n+1}{2n^2-5},
\]
\[
v_n=\frac{\sin n}{\sqrt n},
\]
\[
w_n=\sqrt{n^2+3n}-n.
\]

\subsection*{Exercice 3 -- Suite récurrente}

Soit :
\[
u_0\in[0,1],
\qquad
u_{n+1}=\frac{u_n+1}{2}.
\]

\begin{enumerate}
    \item Montrer que :
    \[
    \forall n\in\N,\quad u_n\in[0,1].
    \]
    \item Étudier la monotonie de \((u_n)\).
    \item Montrer que \((u_n)\) converge.
    \item Déterminer sa limite.
\end{enumerate}

\subsection*{Exercice 4 -- Suites extraites}

Soit \((u_n)\) une suite réelle bornée.  
\begin{enumerate}
    \item Énoncer le théorème de Bolzano-Weierstrass.
    \item Montrer que \((u_n)\) possède une suite extraite convergente.
    \item Donner un exemple de suite bornée divergente.
\end{enumerate}

\subsection*{Exercice 5 -- Équivalents}

Donner un équivalent simple de :
\[
a_n=\frac{5n^4-3n+1}{2n^2+n},
\]
puis de :
\[
b_n=\sqrt{n^2+2n+3}.
\]

\newpage

%=========================================================
\section{Correction détaillée du devoir bilan}

\subsection*{Correction de l'exercice 1}

\textbf{1.}  
On dit que :
\[
u_n\to\ell
\]
si :
\[
\forall \varepsilon>0,\quad \exists N\in\N,\quad \forall n\geq N,\quad \abs{u_n-\ell}\leq\varepsilon.
\]

\textbf{2.}  
Supposons :
\[
u_n\to\ell
\quad\text{et}\quad
u_n\to m.
\]
Soit \(\varepsilon>0\).  
À partir d'un certain rang :
\[
\abs{u_n-\ell}\leq\frac\varepsilon2
\]
et :
\[
\abs{u_n-m}\leq\frac\varepsilon2.
\]
Alors :
\[
\abs{\ell-m}
\leq
\abs{\ell-u_n}+\abs{u_n-m}
\leq
\varepsilon.
\]
Comme ceci est vrai pour tout \(\varepsilon>0\), on obtient :
\[
\ell=m.
\]

\textbf{3.}  
Si \(u_n\to\ell\), alors à partir d'un certain rang :
\[
\abs{u_n-\ell}\leq1.
\]
Donc :
\[
\abs{u_n}\leq\abs{\ell}+1.
\]
Les premiers termes étant en nombre fini, ils sont bornés.  
La suite entière est donc bornée.

\textbf{4.}  
Le théorème des gendarmes affirme que si :
\[
u_n\leq v_n\leq w_n
\]
à partir d'un certain rang, et si :
\[
u_n\to\ell,\qquad w_n\to\ell,
\]
alors :
\[
v_n\to\ell.
\]
La démonstration repose sur l'encadrement :
\[
\ell-\varepsilon\leq u_n\leq v_n\leq w_n\leq\ell+\varepsilon
\]
à partir d'un certain rang.

\textbf{5.}  
Toute suite réelle croissante et majorée converge.  
Toute suite réelle décroissante et minorée converge.

\subsection*{Correction de l'exercice 2}

Pour :
\[
u_n=\frac{3n^2+n+1}{2n^2-5},
\]
on divise par \(n^2\) :
\[
u_n=\frac{3+\frac1n+\frac1{n^2}}{2-\frac5{n^2}}.
\]
Donc :
\[
u_n\to\frac32.
\]

Pour :
\[
v_n=\frac{\sin n}{\sqrt n},
\]
on a :
\[
-\frac1{\sqrt n}\leq v_n\leq\frac1{\sqrt n}.
\]
Les deux bornes tendent vers \(0\), donc :
\[
v_n\to0.
\]

Pour :
\[
w_n=\sqrt{n^2+3n}-n,
\]
on multiplie par la quantité conjuguée :
\[
w_n
=
\frac{(n^2+3n)-n^2}{\sqrt{n^2+3n}+n}
=
\frac{3n}{\sqrt{n^2+3n}+n}.
\]
On factorise \(n\) au dénominateur :
\[
w_n=
\frac{3n}{n\sqrt{1+\frac3n}+n}
=
\frac3{\sqrt{1+\frac3n}+1}.
\]
Donc :
\[
w_n\to\frac3{1+1}=\frac32.
\]

\subsection*{Correction de l'exercice 3}

On considère :
\[
u_{n+1}=\frac{u_n+1}{2}.
\]

\textbf{1.}  
La fonction :
\[
f(x)=\frac{x+1}{2}
\]
envoie \([0,1]\) dans \([0,1]\).  
En effet, si :
\[
0\leq x\leq1,
\]
alors :
\[
1\leq x+1\leq2,
\]
donc :
\[
\frac12\leq f(x)\leq1.
\]
Comme :
\[
u_0\in[0,1],
\]
on obtient par récurrence :
\[
u_n\in[0,1].
\]

\textbf{2.}  
On calcule :
\[
u_{n+1}-u_n
=
\frac{u_n+1}{2}-u_n
=
\frac{1-u_n}{2}.
\]
Comme :
\[
u_n\in[0,1],
\]
on a :
\[
1-u_n\geq0.
\]
Donc :
\[
u_{n+1}-u_n\geq0.
\]
La suite est croissante.

\textbf{3.}  
La suite est croissante et majorée par \(1\).  
Donc elle converge.

\textbf{4.}  
Soit \(\ell\) sa limite.  
Par continuité de \(f\), on obtient :
\[
\ell=\frac{\ell+1}{2}.
\]
Donc :
\[
2\ell=\ell+1,
\]
d'où :
\[
\ell=1.
\]

\subsection*{Correction de l'exercice 4}

\textbf{1.}  
Le théorème de Bolzano-Weierstrass affirme que toute suite réelle bornée possède une suite extraite convergente.

\textbf{2.}  
Comme \((u_n)\) est bornée, le théorème de Bolzano-Weierstrass s'applique directement.  
Donc il existe une extractrice \(\varphi:\N\to\N\) strictement croissante et un réel \(\ell\) tels que :
\[
u_{\varphi(n)}\to\ell.
\]

\textbf{3.}  
La suite :
\[
u_n=(-1)^n
\]
est bornée, car :
\[
\abs{u_n}=1.
\]
Mais elle ne converge pas, puisque :
\[
u_{2n}=1
\]
et :
\[
u_{2n+1}=-1.
\]

\subsection*{Correction de l'exercice 5}

Pour :
\[
a_n=\frac{5n^4-3n+1}{2n^2+n},
\]
le terme dominant du numérateur est :
\[
5n^4.
\]
Le terme dominant du dénominateur est :
\[
2n^2.
\]
Donc :
\[
a_n\sim\frac{5n^4}{2n^2}=\frac52 n^2.
\]

Pour :
\[
b_n=\sqrt{n^2+2n+3},
\]
on écrit :
\[
b_n=n\sqrt{1+\frac2n+\frac3{n^2}}.
\]
Comme :
\[
1+\frac2n+\frac3{n^2}\to1,
\]
on obtient :
\[
\sqrt{1+\frac2n+\frac3{n^2}}\to1.
\]
Donc :
\[
b_n\sim n.
\]

%=========================================================
\newpage
\section{Fiche de synthèse finale}

\subsection*{1. Définition de la limite}

\[
u_n\to\ell
\]
signifie :
\[
\forall \varepsilon>0,\quad \exists N\in\N,\quad \forall n\geq N,\quad \abs{u_n-\ell}\leq\varepsilon.
\]

\subsection*{2. Propriétés fondamentales}

Une limite est unique.

Toute suite convergente est bornée.

Une suite bornée n'est pas forcément convergente.

\subsection*{3. Opérations sur les limites}

Si :
\[
u_n\to\ell,\qquad v_n\to m,
\]
alors :
\[
u_n+v_n\to\ell+m,
\]
\[
u_nv_n\to\ell m.
\]
Si \(m\neq0\), alors :
\[
\frac{u_n}{v_n}\to\frac{\ell}{m}.
\]

\subsection*{4. Inégalités et limites}

Si :
\[
u_n\leq v_n
\]
à partir d'un certain rang, et si :
\[
u_n\to\ell,\qquad v_n\to m,
\]
alors :
\[
\ell\leq m.
\]

\subsection*{5. Théorème des gendarmes}

Si :
\[
u_n\leq v_n\leq w_n
\]
à partir d'un certain rang, et si :
\[
u_n\to\ell,\qquad w_n\to\ell,
\]
alors :
\[
v_n\to\ell.
\]

\subsection*{6. Suites monotones}

Une suite croissante et majorée converge.

Une suite décroissante et minorée converge.

Une suite croissante non majorée tend vers \(+\infty\).

Une suite décroissante non minorée tend vers \(-\infty\).

\subsection*{7. Suites adjacentes}

Deux suites \((u_n)\) et \((v_n)\) sont adjacentes si :
\[
(u_n)\text{ croissante},\qquad (v_n)\text{ décroissante},
\]
\[
u_n\leq v_n,
\]
et :
\[
v_n-u_n\to0.
\]
Alors elles convergent vers la même limite.

\subsection*{8. Suites extraites}

Une suite extraite est une suite :
\[
(u_{\varphi(n)})
\]
où :
\[
\varphi:\N\to\N
\]
est strictement croissante.

Si :
\[
u_n\to\ell,
\]
alors toute suite extraite converge vers \(\ell\).

\subsection*{9. Bolzano-Weierstrass}

Toute suite réelle bornée possède une suite extraite convergente.

\subsection*{10. Suites récurrentes}

Pour :
\[
u_{n+1}=f(u_n),
\]
si \(u_n\to\ell\) et si \(f\) est continue, alors :
\[
f(\ell)=\ell.
\]

\subsection*{11. Comparaisons}

\[
u_n=o(v_n)
\ssi
\frac{u_n}{v_n}\to0.
\]

\[
u_n=O(v_n)
\ssi
\left(\frac{u_n}{v_n}\right)\text{ est bornée à partir d'un certain rang}.
\]

\[
u_n\sim v_n
\ssi
\frac{u_n}{v_n}\to1.
\]

\subsection*{12. Équivalents usuels}

\[
n+1\sim n.
\]

Si \(P\) est un polynôme de degré \(p\) et de coefficient dominant \(a_p\), alors :
\[
P(n)\sim a_p n^p.
\]

Si :
\[
u_n\sim v_n
\]
et :
\[
v_n\to\ell,
\]
alors :
\[
u_n\to\ell.
\]

%=========================================================
\section*{Conclusion pédagogique}
\addcontentsline{toc}{section}{Conclusion pédagogique}

Les suites numériques constituent le premier outil fondamental de l'analyse.  
Elles permettent de comprendre rigoureusement la notion de limite et préparent directement :
\begin{itemize}
    \item l'étude des fonctions ;
    \item les séries numériques ;
    \item les suites de fonctions ;
    \item les intégrales généralisées ;
    \item la topologie des espaces vectoriels normés.
\end{itemize}

Les méthodes essentielles à retenir sont :
\begin{itemize}
    \item revenir à la définition avec \(\varepsilon\) ;
    \item utiliser les opérations sur les limites ;
    \item encadrer avec le théorème des gendarmes ;
    \item exploiter la monotonie et les bornes ;
    \item extraire des sous-suites pour prouver une divergence ;
    \item utiliser les équivalents pour simplifier les calculs asymptotiques.
\end{itemize}

\begin{center}
\[
\boxed{
\text{Étudier une suite, c'est comprendre son comportement à l'infini.}
}
\]
\end{center}

\end{document}